% SOURCE_AUTHORITY: sga5_fr_workpass.tex, lines 7578--7632 (LF-normalized unit).
% SOURCE_SHA256: 791F4EFFC5E02832D5D77ED03518C8156D6F07E4C8238B03545DB93D883FBB28
\subsection*{5. Sistemas projetivos \(J\)-ádicos e AR-\(J\)-ádicos noetherianos}

Em todo este número, são dados um anel comutativo noetheriano \(A\), um ideal \(J\) de \(A\) e uma \(A\)-categoria abeliana \(\mathcal C\).

\subsubsection*{5.1. Generalidades}

\begin{statement}{Definição 5.1.1}
Um objeto \(X\) da categoria \(\Hom(\mathbb N^\circ,\mathcal C)\) diz-se \(J\)-ádico noetheriano (resp. artiniano) se é \(J\)-ádico e seus componentes são objetos noetherianos (resp. artinianos) da categoria \(\mathcal C\).
\end{statement}

Denotar-se-á por \(J\)-adn\((\mathcal C)\) (resp. \(J\)-adf\((\mathcal C)\)) a subcategoria plena de
\(P=\Hom(\mathbb N^\circ,\mathcal C)\) cujos objetos são os sistemas projetivos \(J\)-ádicos noetherianos (resp. artinianos). É claro (3.1.3) que uma extensão de sistemas projetivos \(J\)-ádicos noetherianos (resp. artinianos) é do mesmo tipo.

\begin{statement}{Proposição 5.1.2}
Seja \(X\) um objeto \(J\)-ádico de \(P\). Para que \(X\) seja \(J\)-ádico noetheriano (resp. artiniano), é necessário e suficiente que seja \(J\)-ádico e que \(X_0\) seja noetheriano (resp. artiniano).
\end{statement}

\begin{prooftext}
Seja, com efeito, \(X\) um sistema projetivo \(J\)-ádico tal que \(X_0\) seja noetheriano (resp. artiniano). Consideremos então o epimorfismo canônico (4.2.4):
\[
\gr_J(A)\otimes_A X_0\longrightarrow \grs(X).
\]
Como \(X_0\) é noetheriano (resp. artiniano) e \(J^n\) é de tipo finito para todo \(n\), a existência deste epimorfismo mostra que, para todo inteiro \(n\), o núcleo do morfismo canônico
\[
X_n\longrightarrow X_{n-1}
\]
É noetheriano (resp. artiniano). Deduz-se imediatamente por indução sobre \(n\) que os \(X_n\) são noetherianos (resp. artinianos).
\end{prooftext}

\begin{statement}{Definição 5.1.3}
Seja \(X\) um sistema projetivo indexado por \(\mathbb N\) com valores em \(\mathcal C\). Diz-se que é AR-\(J\)-ádico noetheriano (resp. artiniano) se é AR-\(J\)-ádico e seus componentes são objetos noetherianos (resp. artinianos).
\end{statement}

Denotar-se-á por AR-\(J\)-adn\((\mathcal C)\) (resp. AR-\(J\)-adf\((\mathcal C)\)) a subcategoria plena de \(\underline{\Hom}_{\mathrm{AR}}(\mathbb N^\circ,\mathcal C)\) (2.4) engendrada pelas imagens dos sistemas projetivos AR-\(J\)-ádicos noetherianos (resp. artinianos).

Então é claro que os funtores canônicos
\[
p^n_{\mathrm{AR}}:J\text{-adn}(\mathcal C)\longrightarrow \mathrm{AR}\text{-}J\text{-adn}(\mathcal C),
\]
\[
p^f_{\mathrm{AR}}:J\text{-adf}(\mathcal C)\longrightarrow \mathrm{AR}\text{-}J\text{-adf}(\mathcal C),
\]
obtidos por restrição de \(p_{\mathrm{AR}}\) (2.4) são equivalências de categorias.

Recordemos brevemente o enunciado e a demonstração do teorema de Hilbert.

\begin{statement}{Teorema 5.1.4}
Teorema de Hilbert. Sejam \(A\) um anel comutativo, \(\mathcal C\) uma \(A\)-categoria abeliana e \(M\) um objeto de tipo finito (resp. noetheriano) de \(\mathcal C\). Para toda \(A\)-álgebra de tipo finito \(B\), graduada por \(\mathbb N\), o \((B,\mathcal C)\)-módulo graduado
\[
B\otimes_A M
\]
É de tipo finito (resp. noetheriano).

N.B. Se \(B=\bigoplus_{n\ge0}B_n\), \(B\otimes_A M\) designa o objeto graduado \(\bigl(B_n\otimes_A M\bigr)_{n\in\mathbb N}\).
\end{statement}
